For a Hamiltonian graph \(G\), we study the minimum number \(\hsep(G)\) of
Hamiltonian cycles needed to distinguish every ordered pair of distinct edges:
for each \((e,f)\), some selected cycle must contain \(e\) and avoid \(f\).
We report an exhaustively certified finite computation over all 22,263
non-isomorphic Barnette graphs of orders at most 36.  The resulting extremal
values are independently checked by explicit Hamiltonian-cycle covers and
matching lower certificates or complete-universe exclusion.  In particular,
the maximum at order 36 is 49 and is attained uniquely by the graph with
canonical hash
{\scriptsize\nolinkurl{2f96ada16c46cd2bd038b97d5af44f46ed14522cba1edf3fa041d66e02107bcc}}.
The unique maximizers at orders divisible by four form, over the certified
range, a double-ladder family linked by a fixed four-vertex square insertion.
We define this family, prove its Hamiltonian-cycle classification, and prove
for all odd \(a,b\ge3\) the exact formula
\[
 \hsep(D(a,b))=\frac{(a+2)(b+2)-1}{2}.
\]
The lower bound is a packing of ordered edge requirements whose supports are
exception singletons, boundary-turn singletons, and disjoint interior
dominoes.  The matching upper bound selects the four connector-exception
cycles, every boundary turn, and one parity class of interior turns; the
resulting edge-incidence signatures form an antichain.  Three values frozen
before computation and later certified for \(D(9,9)\), \(D(11,9)\), and
\(D(11,11)\) provide prospective finite confirmation.  They do not prove
extremality at orders 40, 44, or 48: general extremality of balanced double
ladders remains open.  The terminology and its relationship to prior
separation notions also remain subject to a systematic literature review.
